\documentclass{ug}
\usepackage{xltabular}

%replace ipcore-name by the name of your ip core (e.g. IOb-Cache) and description by a brief description (e.g. a Configurable Cache)

\title{%
\Huge IOb-SoC \\
 \vspace*{3cm}
\Large System-on-Chip (SoC) template 
}

\header{IOb-SoC, System-on-Chip (SoC) template}

\date{\today}
\category{User Guide, \input{\NAME_version.tex}, Build \input{shortHash.tex}}

\definecolor{iob-green}{rgb}{0.0,1.0,0.80}
\definecolor{iob-blue}{rgb}{0.90196,0.94902,1}
\definecolor{dark-blue}{RGB}{155,194,230}
\definecolor{mid-blue}{RGB}{189,215,238}
\definecolor{light-blue}{RGB}{221,235,247}


\begin{document}

\maketitle
\pagenumbering{gobble}

\vspace*{\fill}
User Guide, \input{\NAME_version.tex}, Build \input{shortHash.tex}
\hspace*{\fill} \includegraphics[keepaspectratio,scale=.7]{Logo.png}

\cleardoublepage
\pagenumbering{roman}
\setcounter{page}{1}
\today & Initial document version \input{version}. \\ \hline

\cleardoublepage
\tableofcontents
\clearpage
\listoftables
\clearpage
\listoffigures
\cleardoublepage
\pagenumbering{arabic}
\setcounter{page}{1}
\section{Introduction}
\label{sec:intro}
\lipsum[1-1]

\begin{figure}[!htbp]
  \centerline{\includegraphics[width=14cm]{symb.pdf}}
  \vspace{0cm}
  \caption{IP Core Symbol}
  \label{fig:symbol}
\end{figure}


\subsection{Features}
\label{sec:feat}
\begin{itemize}
\item 32-bit RISC-V control CPU
\item Support for Integer (I), atomic (A) and multiply/divide extensions (M)
\item Instruction and data caches
\item RS232 interfaces for viewing runtime messages
\item Optional timer peripheral
\item Optional Ethernet peripheral
\item Frequency of operation at 167MHz on FPGA
\item Needs external DDR4 memory controller IP
\end{itemize}


\subsection{Benefits}
\label{sec:benef}
\begin{itemize}
  \itemsep-0.5em
\item Compact and easy to integrate hardware and software implementation
\item Low power consumption
\end{itemize}


\subsection{Deliverables}
\label{sec:deliv}
\begin{itemize}
  \itemsep-0.5em
\item FPGA synthesized netlist or Verilog source code, and respective
  synthesis and implementation scripts
\item FPGA verification environment by simulation and emulation
\item Bare-metal software driver and example user software
\item User documentation for easy system integration
\end{itemize}


\section{Description}

This section gives a detailed description of the IP core. The high-level block
diagram is presented, along with a description of its blocks. The parameters and
macros that define the core configuration are listed and explained. The
interface signals are enumerated and described; if timing diagrams are needed,
they are shown after the interface signals. Finally the Control and Status
Registers (CSR) are outlined and explained.

\subsection{Block Diagram}
\label{sec:bdd}
Figure~\ref{fig:bd} presents a high-level block diagram of the core, followed by
a brief description of each block.

\begin{figure}[H]
  \centering {\includegraphics[width=\columnwidth]{bd.pdf}}
  \vspace{-0.7cm}
  \caption{High-Level Block Diagram}
  \label{fig:bd}
\end{figure}


\setlist[description]{font=\bfseries\MakeUppercase, labelindent=\parindent, leftmargin=*}

\input{blocks}


\ifdefined\CONFS
\subsection{Configuration}
\label{sec:ipconfig}
Table~\ref{tab:confs} describes the IP core configuration. The core may be configured using macros or parameters:

\begin{description}
\item \textbf{'M'} Macro: a Verilog macro or ``\`define'' directive is used to include or exclude code segments, to create core configurations that are valid for all instances of the core.
\item \textbf{'P'} Parameter: a Verilog parameter is passed to each instance of the core and defines the configuration of that particular instance.
\end{description}

\begin{xltabular}{\textwidth}{|l|c|c|c|c|X|} \hline
    \rowcolor{iob-green}
    {\bf Configuration} & {\bf Type} & {\bf Min} & {\bf Typical} & {\bf Max} & {\bf Description}
    \\ \hline \hline
    \input confs_tab
    \caption{Core Configuration.}\label{tab:confs}
\end{xltabular}

The parameters in the top-level Verilog module that are not listed above are
called Derived Parameters. They are given as function of the primary parameters
and should never be changed. They are used to simplify the definition of the
interface and internal signals. The list of derived parameters is given below:
\input derived_params

\fi

\subsection{Interface Signals}
\label{sec:ifsig}
\begin{table}[H]
  \centering
  \begin{tabular}{|l|l|r|p{10.5cm}|}

    \hline
    \rowcolor{iob-green}
    {\bf Name} & {\bf Direction} & {\bf Width} & {\bf Description}  \\ \hline \hline

    \input{gen_if_tab}

  \end{tabular}
  \caption{General Interface Signals.}
  \label{gen_if_tab:is}
\end{table}

\begin{table}[H]
  \centering
  \begin{tabular}{|l|l|r|p{8.5cm}|}

    \hline
    \rowcolor{iob-green}
    {\bf Name} & {\bf Direction} & {\bf Width} & {\bf Description}  \\ \hline \hline

    \input{iob_s_if_tab}

  \end{tabular}
  \caption{IObundle Interface Signals}
  \label{tab:if_iob_s}
\end{table}

\begin{table}[H]
  \centering
  \begin{tabular}{|l|l|r|p{9.5cm}|}
    
    \hline
    \rowcolor{iob-green}
    {\bf Name} & {\bf Direction} & {\bf Width} & {\bf Description}  \\ \hline \hline

    \input{rs232_if_tab}
 
  \end{tabular}
  \caption{RS232 Interface Signals}
  \label{tab:if_rs232}
\end{table}

%TODO
%\input{\TEX/ug/axil_s_if}


%timing diagrams
\ifdefined\TD
\subsection{Timing Diagrams}
\label{sec:td}
\input{td}
\fi

%software components
\ifdefined\SWREG
\subsection{Control and Status Registers}
\label{sec:swreg}
\begin{table}[H]
  \centering
  \begin{tabular}{|l|c|c|r|P{0.8cm}|p{5.5cm}|}
    \hline
    \rowcolor{iob-green}
    
    {\bf Name} & {\bf R/W} & {\bf Addr} & {\bf Bits} & {\bf Initial Value} & {\bf Description} \\ \hline

    \input{uart_swreg_tab}
    
  \end{tabular}
  \caption{UART software accessible registers.}
  \label{tab:sw_tx}
\end{table}


\fi

\section{Usage}

\subsection{Instantiation}
\label{sec:inst}
Figure~\ref{fig:inst} illustrates how to instantiate the IP core and, if applicable, the
required external blocks.

\begin{figure}[!htbp]
    \centerline{\includegraphics[width=14cm]{inst.pdf}}
    \vspace{0cm}\caption{Core Instance and Required Surrounding Blocks}
    \label{fig:inst}
\end{figure}

Add here any specific ip core description.




\subsection{Simulation}
\label{sec:tbbd}
The provided testbench uses the core instance described in
Section~\ref{sec:inst}. A high-level block diagram of the testbench is shown in
Figure~\ref{fig:tbbd}. The testbench is organized in a modular fashion, with
each test described in a separate file. The test suite consists of all the test
case files to make adding, modifying, or removing tests easy.

\begin{figure}[!htbp]
    \centerline{\includegraphics[width=14cm]{tbbd.pdf}}
    \vspace{0cm}\caption{Testbench Block Diagram}
    \label{fig:tbbd}
\end{figure}

Add here any specific ip core description.


\subsubsection*{System-level Simulation}

Upon request, simulation files to run the core embedded in a RISC-V system can
be provided. The core is exercised in various modes by the RISC-V processor,
using a bare-metal software program written in the C programming language.


\ifdefined\ASICSYNTH
\subsection{ASIC Synthesis}
\label{sec:synth}
This IP core synthesizes for FPGA and ASIC with very low logic resources
consumption. The implementation does not require memory or arithmetic resources.

\fi

\ifdefined\FPGACOMP
\subsection{FPGA Compilation}
\label{sec:fpga}
The Quartus and Vivado FPGA compilation toolchains are supported via {\tt .tcl}
scripts invoked by a Makefile. The script compiles and elaborates the design for
a given set of target FPGA boards. Additional boards can be supported by
following the flow provided.

The core is instantiated within a Verilog wrapper, which depends on the
specifics of the FPGA device and the FPGA board. Timing constraints file(s) are
provided.

After compilation, reports on FPGA resource usage, power consumption, and timing
closure are generated. A post-synthesis Verilog file is created, which can be
used in post-synthesis simulation.

Add here any specific ip core description.


\subsubsection*{System-level FPGA Run}

Upon request, files to run the core on the supported FPGA boards can be
provided. The core is embedded in a RISC-V system and exercised in various
modes, using a bare-metal software program written in the C programming
language.

\fi

\ifdefined\RESULTS
\section{Implementation Results}
\label{sec:results}
The results obtained use all the default user configurable parameters, which affect the RAM usage results.

\fi

\ifdefined\CUSTOM
\input{custom}
\fi

\end{document}
